% !TeX root = ../main.tex

% ==================================================
% 2.6. __new__와 __init__ — 인스턴스 생성과 초기화는 어떻게 다른가?
% ==================================================

\subsection{\zsh{__new__}와 \zsh{__init__} - 인스턴스 생성과 초기화는 어떻게 다른가?}

2.5에서는 클래스에 정의된 함수 객체가 인스턴스에 바인딩되면서 메서드가 호출되는 과정을 학습하였다.

특히 다음 두 호출이 일반적인 인스턴스 메서드에서 같은 효과를 가진다는 사실을 확인하였다.

\begin{dhkcode}
p.introduce()

Person.introduce(p)
\end{dhkcode}

이번 2.6에서는 그보다 앞선 과정, 즉 인스턴스가 존재하기 전부터 인스턴스가 생성되고 초기화될 때까지 Python에서 어떤 일이 발생하는지 살펴볼 것이다.

출발점은 다음 코드다.

\begin{dhkcode}
class Person:
    def __init__(self, name):
        self.name = name


p = Person("Kim")
\end{dhkcode}

우리는 흔히 \zsh{__init__()}을 생성자(Constructor)라고 부르지만,

엄밀하게는 \zsh{__init__()}이 인스턴스 자체를 생성하는 것은 아니다.

Python에서는 일반적으로 다음 두 과정을 구분한다.

\begin{dhksmry}
    __new__() : 인스턴스 생성

    __init__() : 생성된 인스턴스 초기화
\end{dhksmry}

이번 학습의 핵심은 이 두 메서드가 어떤 순서로 호출되고, 각각 무엇을 반환하며, 그 과정에서 \zsh{self}가 어떻게 전달되는지를 이해하는 것이다.

\subsubsection{개념 정의 - 생성과 초기화의 차이}

먼저 객체의 생성(Creation)과 초기화(Initialization)를 구분하자.

\begin{dhknote}
\begin{dhksmry}[centered]
인스턴스 생성
\end{dhksmry}
\end{dhknote}

아직 존재하지 않는 인스턴스 객체를 만들어 반환하는 과정이다.

\begin{dhknote}
\begin{dhksmry}[centered]
인스턴스 초기화
\end{dhksmry}
\end{dhknote}

이미 생성된 인스턴스에 필요한 속성을 설정하거나 초기 상태를 구성하는 과정이다.

예를 들어 다음 코드를 보자.

\begin{dhkcode}
class Person:
    def __init__(self, name):
        self.name = name
        self.age = 0
\end{dhkcode}

여기서 \zsh{__init__()}은 인스턴스에 \zsh{name}과 \zsh{age} 속성을 설정하고 있다.

하지만 \zsh{self.name}과 \zsh{self.age}를 설정하려면 그 속성들을 담을 인스턴스 객체가 먼저 존재해야 하지 않을까?

바로 이 질문이 \zsh{__new__()}와 \zsh{__init__()}을 구분하는 출발점이다.

\begin{dhktable}
|구분|	__new__()|	__init__()|
|:---|:---|:---|
|역할|	인스턴스 생성|	인스턴스 초기화|
|첫 번째 매개변수 관례|	cls|	self|
|첫 번째 인수|	생성 대상 클래스|	생성된 인스턴스|
|반환값|	생성한 객체 또는 다른 객체|	반드시 None|
|실행 순서|	일반적으로 먼저|	생성 이후|
|직접 재정의하는 빈도|	상대적으로 드묾|	매우 흔함|
\end{dhktable}

여기서 \zsh{cls}와 \zsh{self}는 모두 관례적인 매개변수 이름이다.

\zsh{self}가 예약어가 아니었던 것처럼 \zsh{cls}도 예약어는 아니다.

\subsubsection{클래스 객체를 호출하면 무엇이 실행되는가?}

다음 표현식을 살펴보자.

\begin{dhkcode}
Person("Kim")
\end{dhkcode}

2.3에서는 이 표현식이 클래스 객체를 호출하는 것이라고 학습했다.

그렇다면 클래스 객체를 호출할 때 Python은 무엇을 실행할까?

일반적인 Python 클래스는 \zsh{type}이라는 메타클래스의 인스턴스다.

\begin{dhkcode}
class Person:
    pass

print(type(Person))
\end{dhkcode}

실행 결과:

\begin{dhkzsh}
<class 'type'>
\end{dhkzsh}

여기서 중요한 사실이 있다.

클래스 객체도 호출 가능한 객체이며, 그 호출 과정은 메타클래스의 호출 메커니즘에 의해 처리된다는 것이다.

일반적인 클래스에서는 \zsh{type.__call__()}이 인스턴스 생성 과정을 담당한다.

개념적인 실행 순서는 다음과 같다.

\begin{dhkgraph}[uniform]
node n1 (1,1) :: Person('Kim') // 클래스 객체 호출
node n2 (2,1) :: type.__call__ // 생성 과정 조정
node n3 (3,1) [accent] :: Person.__new__ // 인스턴스 생성 시도
node n4 (4,1) [accent] :: Person.__init__ // 생성된 인스턴스 초기화
node n5 (5,1) :: 호출 결과 반환 // p 이름에 바인딩

edge n1->n2
edge n2->n3
edge n3->n4
edge n4->n5

caption :: 기본 메타클래스 type을 사용하는 일반적인 클래스의 인스턴스 생성 과정
\end{dhkgraph}

단, 이 그림에는 한 가지 조건이 있다.

\zsh{__new__()}가 해당 클래스의 인스턴스 객체를 반환한 경우에만 \zsh{__init__()}이 이어서 호출된다.

다른 종류의 객체를 반환하면 \zsh{__init__()}이 호출되지 않을 수 있다.

이 예외는 뒤에서 실제 코드로 확인할 것이다.

\begin{dhknote}
\begin{dhksmry}[centered]
type.__call__()을 개념적으로 표현하면
\end{dhksmry}
\end{dhknote}

다음과 비슷하게 생각할 수 있다.

\begin{dhkcode}
def conceptual_call(cls, *args, **kwargs):
    obj = cls.__new__(cls, *args, **kwargs)

    if isinstance(obj, cls):
        cls.__init__(obj, *args, **kwargs)

    return obj
\end{dhkcode}

이 코드는 실제 CPython 구현을 그대로 옮긴 것이 아니라, 일반적인 동작 순서를 이해하기 위한 수도 코드다.

실제 구현에서는 특수 메서드 탐색, 타입 검사, 초기화 결과 검증 등의 세부 처리가 추가된다.

또한 위 수도 코드는 \zsh{__init__()}의 반환값이 \zsh{None}이어야 한다는 검사도 생략하였다.

현재 단계에서 중요한 것은 다음 순서다.

\begin{dhksmry}[centered]
Class Call \ra __new__ \ra __init__ \ra Return
\end{dhksmry}

\subsubsection{\zsh{__new__()} - 인스턴스 생성}

이제 \zsh{__new__()}를 직접 정의해 보자.

\begin{dhkcode}
class Person:
    def __new__(cls):
        print("__new__ 실행")

        instance = super().__new__(cls)

        return instance
\end{dhkcode}

이 코드에서 가장 중요한 부분은 다음 표현식이다.

\begin{dhkcode}
instance = super().__new__(cls)
\end{dhkcode}

여기서는 상위 클래스에서 제공하는 인스턴스 생성 기능을 호출한다.

\zsh{Person}의 직접적인 상위 클래스는 기본적으로 \zsh{object}이므로,

이 예제에서는 \zsh{object.__new__()}를 통해 인스턴스를 생성하는 효과가 있다.

\zsh{super()}의 정확한 메서드 탐색 규칙은 2.9에서 다룰 것이다.

\begin{dhknote}
\begin{dhksmry}[centered]
cls는 무엇을 참조할까?
\end{dhksmry}
\end{dhknote}

다음 코드로 확인해 보자.

\begin{dhkcode}
class Person:
    def __new__(cls):
        print(cls)

        return super().__new__(cls)

p = Person()
\end{dhkcode}

실행 결과:

\begin{dhkzsh}
<class '__main__.Person'>
\end{dhkzsh}

즉, \zsh{cls}는 생성 대상인 \zsh{Person} 클래스 객체를 참조한다.

2.5에서 \zsh{self}가 인스턴스 객체를 참조했던 것과 비교해 보자.

\begin{dhktable}
|매개변수|참조하는 객체|
|:---|:---|
|cls|클래스 객체|
|self|인스턴스 객체|
\end{dhktable}

이 차이는 자연스러운 결과다.

\zsh{__new__()}이 호출되는 시점에는 새 인스턴스를 생성하려는 중이므로, 아직 그 인스턴스를 \zsh{self}로 전달할 수 없기 때문이다.

\begin{dhknote}
\begin{dhksmry}[centered]
생성한 인스턴스를 확인해 보자.
\end{dhksmry}
\end{dhknote}

\begin{dhkcode}
class Person:
    def __new__(cls):
        instance = super().__new__(cls)

        print(type(instance))

        return instance

p = Person()
\end{dhkcode}

실행 결과:

\begin{dhkzsh}
<class '__main__.Person'>
\end{dhkzsh}

\zsh{super().__new__(cls)}가 반환한 객체는 \zsh{Person} 클래스의 인스턴스다.

그리고 \zsh{__new__()}는 그 객체를 반환하고 있다.

여기서 중요한 사실은 \zsh{__new__()}가 반드시 새 인스턴스를 반환해야 하는 것은 아니라는 것이다.

다른 객체를 반환하는 것도 가능하다.

다만 그렇게 하면 이후 초기화 과정이 달라질 수 있다.

\subsubsection{\zsh{__init__()} - 생성된 인스턴스 초기화}

이번에는 \zsh{__init__()}을 살펴보자.

\begin{dhkcode}
class Person:
    def __init__(self, name):
        print("__init__ 실행")

        self.name = name

p = Person("Kim")
\end{dhkcode}

실행 결과:

\begin{dhkzsh}
__init__ 실행
\end{dhkzsh}

여기서 \zsh{self}는 이미 생성된 \zsh{Person} 인스턴스를 참조한다.

따라서 다음 코드를 실행할 수 있다.

\begin{dhkcode}
self.name = name
\end{dhkcode}

2.4에서 배운 것처럼, 이 코드는 \zsh{self}가 참조하는 인스턴스에 \zsh{name} 속성을 설정하는 동작이다.

실제로 확인해 보자.

\begin{dhkcode}
print(p.__dict__)
\end{dhkcode}

실행 결과:

\begin{dhkzsh}
{'name': 'Kim'}
\end{dhkzsh}

\begin{dhknote}
\begin{dhksmry}[centered]
__init__()은 객체를 반환하는가?
\end{dhksmry}
\end{dhknote}

여기서 주의할 점이 있다.

\zsh{__init__()}은 생성된 인스턴스를 초기화하는 메서드이므로 반드시 \zsh{None}을 반환해야 한다.

다음 코드를 보자.

\begin{dhkcode}
class Person:
    def __init__(self):
        return 10


p = Person()
\end{dhkcode}

실행하면 다음과 같은 오류가 발생한다.

\begin{dhkzsh}
TypeError: __init__() should return None, not 'int'
\end{dhkzsh}

왜 그럴까?

\zsh{__init__()}은 인스턴스 생성 결과를 결정하는 메서드가 아니라, 이미 생성된 인스턴스를 초기화하는 메서드이기 때문이다.

일반적인 함수에서는 반환값을 자유롭게 지정할 수 있지만, Python의 특수 메서드에는 각각의 동작 규약이 존재한다.

\zsh{__init__()}도 그 규약을 따르는 것이다.

다음처럼 명시적으로 \zsh{None}을 반환하는 것은 가능하다.

\begin{dhkcode}
class Person:
    def __init__(self):
        return None
\end{dhkcode}

하지만 보통 \zsh{return}을 작성하지 않는다.

Python 함수는 명시적인 반환문 없이 종료되면 자동으로 \zsh{None}을 반환하기 때문이다.

\subsubsection{\zsh{__new__()}와 \zsh{__init__()}의 실행 순서 확인}

이제 두 메서드를 함께 정의해 보자.

\begin{dhkcode}
class Person:
    def __new__(cls, name):
        print("1. __new__ 실행")

        instance = super().__new__(cls)

        print("2. 인스턴스 생성 완료")

        return instance

    def __init__(self, name):
        print("3. __init__ 실행")

        self.name = name

        print("4. 인스턴스 초기화 완료")


p = Person("Kim")

print("5. p에 바인딩된 인스턴스:", p.name)
\end{dhkcode}

실행 결과:

\begin{dhkzsh}
1. __new__ 실행
2. 인스턴스 생성 완료
3. __init__ 실행
4. 인스턴스 초기화 완료
5. p에 바인딩된 인스턴스: Kim
\end{dhkzsh}

여기서 주목할 점은 \zsh{__new__()}와 \zsh{__init__()} 모두 \zsh{"Kim"}을 전달받았다는 사실이다.

\begin{dhkcode}
def __new__(cls, name):

def __init__(self, name):
\end{dhkcode}

하지만 첫 번째 매개변수가 참조하는 객체는 서로 다르다.

\begin{dhkbranch}[in=8]
node main :: Person("Kim")
arrow main
node main :: __new__(cls, name)
branchline main -> cls :: cls -> Person 클래스 객체
stop cls
branchline main -> name :: name -> "Kim"
stop name
arrow main
node main :: 새 인스턴스 생성
arrow main
node main :: __init__(self, name)
branchline main -> self :: self -> 생성된 Person 인스턴스
stop self
branchline main -> nn :: name -> "Kim"
stop nn
arrow main
node main :: self.name = name
arrow main
node main :: 초기화 완료
arrow main
node main :: p에 인스턴스 바인딩
\end{dhkbranch}

이 구조는 이전에 학습했던 함수의 인수 전달 및 이름 바인딩 규칙과 연결된다.

\zsh{cls, self, name}은 모두 함수 실행 과정에서 객체를 참조하도록 바인딩되는 매개변수 이름이다.

차이는 Python이 특수 메서드를 호출할 때 어떤 객체를 전달하도록 정해 놓았는가에 있다.

\subsubsection{\zsh{__new__()}가 다른 객체를 반환하면 어떻게 되는가?}

앞에서 \zsh{__new__()}는 반드시 새로운 인스턴스를 반환할 필요가 없다고 했다.

이번에는 그 의미를 실제 코드로 확인해 보자.

\begin{dhkcode}
class Person:
    def __new__(cls):
        print("__new__ 실행")

        return "Not a Person"

    def __init__(self):
        print("__init__ 실행")


p = Person()

print(p)
print(type(p))
\end{dhkcode}

실행 결과:

\begin{dhkzsh}
__new__ 실행
Not a Person
<class 'str'>
\end{dhkzsh}

여기서 중요한 사실이 있다.

\zsh{__init__ 실행} 이라는 메시지가 출력되지 않았다.

왜 그럴까?

\zsh{__new__()}가 \zsh{Person} 인스턴스가 아니라 문자열 객체를 반환했기 때문이다.

일반적인 \zsh{type.__call__()} 동작에서는 \zsh{__new__()}의 반환값이 생성 대상 클래스의 인스턴스가 아니면 \zsh{__init__()}을 호출하지 않는다.

따라서 이 예제에서는 다음과 같은 과정이 이루어진다.

\begin{dhksmry}[centered]
Person()
\da
Person.__new__()
\da
문자열 객체 반환
\da
Person 인스턴스가 아님
\da
Person.__init__() 호출하지 않음
\da
문자열 객체 반환
\end{dhksmry}

이 예제를 통해 중요한 사실을 알 수 있다.

클래스 객체를 호출한다고 해서 반드시 그 클래스의 인스턴스가 반환되는 것은 아니다.

일반적으로는 인스턴스를 반환하도록 설계하지만, 특수한 클래스에서는 다른 동작을 구현할 수 있다.

이 점은 클래스 호출이 단순한 객체 복제 문법이 아니라, 정해진 호출 프로토콜을 통해 실행되는 과정이라는 사실을 보여준다.

\subsubsection{\zsh{__new__()}를 직접 재정의하는 이유}

그렇다면 왜 \zsh{__new__()}와 \zsh{__init__()}을 분리했을까?

일반적인 클래스에서는 \zsh{__init__()}만 정의해도 충분한 경우가 많다.

\begin{dhkcode}
class Person:
    def __init__(self, name):
        self.name = name
\end{dhkcode}

그런데 다음과 같은 상황에서는 인스턴스 생성 과정 자체를 제어해야 할 수 있다.

\begin{dhksmry}
    1. 생성할 객체의 타입을 결정해야 하는 경우

    2. 불변 객체의 값을 생성 시점에 설정해야 하는 경우

    3. 이미 존재하는 객체를 반환해야 하는 경우

    4. 인스턴스 생성을 제한하거나 관리해야 하는 경우
\end{dhksmry}

대표적인 예로 불변 객체를 생각해 볼 수 있다.

\begin{dhknote}
\begin{dhksmry}[centered]
불변 객체를 상속하는 경우
\end{dhksmry}
\end{dhknote}

다음 코드를 보자.

\begin{dhkcode}
class PositiveInt(int):
    def __new__(cls, value):
        if value < 0:
            raise ValueError("양수 또는 0만 허용됩니다.")

    return super().__new__(cls, value)
\end{dhkcode}

이제 다음처럼 사용할 수 있다.

\begin{dhkcode}
a = PositiveInt(10)

print(a)
print(type(a))
\end{dhkcode}

실행 결과:

\begin{dhkzsh}
10
<class '__main__.PositiveInt'>
\end{dhkzsh}

반면 다음 코드는 오류를 발생시킨다.

\begin{dhkcode}
b = PositiveInt(-5)
\end{dhkcode}

\begin{dhkzsh}
ValueError: 양수 또는 0만 허용합니다.
\end{dhkzsh}

이 예제에서 \zsh{PositiveInt}는 \zsh{int}를 상속받아 새로운 정수형 객체를 정의했다.

\zsh{int}는 불변 객체이므로 정수값 자체는 생성 과정에서 결정된다.

따라서 이런 경우에는 \zsh{__init__()} 보다 \zsh{__new__()}를 재정의하는 것이 적절하다.

이처럼 \zsh{__new__()}는 인스턴스가 만들어지는 과정 자체를 제어해야 할 때 사용된다.

\subsubsection{\zsh{__new__()}와 \zsh{__init__()}을 직접 호출하면 어떻게 되는가?}

이번에는 두 메서드를 직접 호출해 보자.

먼저 다음 클래스를 정의하자.

\begin{dhkcode}
class Person:
    def __init__(self, name):
        self.name = name
\end{dhkcode}

\begin{dhknote}
\begin{dhksmry}[centered]
__new__() 직접 호출
\end{dhksmry}
\end{dhknote}

\begin{dhkcode}
p = Person.__new__(Person)
\end{dhkcode}

이 코드는 \zsh{Person}의 인스턴스를 생성한다.

하지만 \zsh{__init__()}은 호출하지 않았다.

따라서 다음 코드를 실행하면

\begin{dhkcode}
print(p.__dict__)
\end{dhkcode}

실행 결과:

\begin{dhkzsh}
{}
\end{dhkzsh}

인스턴스는 존재하지만 \zsh{name} 속성은 설정되지 않았다.

\begin{dhknote}
\begin{dhksmry}[centered]
__init__() 직접 호출
\end{dhksmry}
\end{dhknote}

이제 다음 코드를 실행해 보자.

\begin{dhkcode}
Person.__init__(p, "Kim")
\end{dhkcode}

그다음 속성을 확인하면

\begin{dhkcode}
print(p.__dict__)
\end{dhkcode}

실행 결과:

\begin{dhkzsh}
{'name': 'Kim'}
\end{dhkzsh}

즉, 직접 호출한 \zsh{__init__()}이 기존 인스턴스를 초기화한 것이다.

이 과정을 정리하면 다음과 같다.

\begin{dhkcode}
p = Person.__new__(Person)

Person.__init__(p, "Kim")
\end{dhkcode}

이는 일반적인 \zsh{Person("Kim")} 호출에서 일어나는 생성 및 초기화 과정을 이해하기 위한 실험이다.

다만 모든 클래스에 대해 이렇게 수동으로 생성과 초기화를 분리해도 안전한 것은 아니다.

일부 클래스는 생성 과정에서 특별한 불변 조건이나 내부 상태를 설정하기 때문에, 일반적으로는 클래스 객체를 직접 호출하는 방식을 사용해야 한다.

\begin{dhkcode}
p = Person("Kim")
\end{dhkcode}

\begin{dhknote}
\begin{dhksmry}[centered]
왜 이 실험이 중요할까?
\end{dhksmry}
\end{dhknote}

이전에는 다음 코드만 봤다.

\begin{dhkcode}
p = Person("Kim")
\end{dhkcode}

이제는 이 한 줄 안에서 개념적으로 두 가지 다른 동작이 이루어진다는 사실을 확인한 것이다.

\begin{dhksmry}[centered]
Instance Creation $\neq$ Instance Initialization
\end{dhksmry}

\subsubsection{\zsh{__init__()}은 인스턴스마다 실행되는가?}

다음 코드를 살펴보자.

\begin{dhkcode}
class Person:
    def __init__(self, name):
        print("__init__ 실행", name)

        self.name = name

p1 = Person("Kim")
p2 = Person("Lee")
\end{dhkcode}

실행 결과:

\begin{dhkzsh}
__init__ 실행: Kim
__init__ 실행: Lee
\end{dhkzsh}

일반적인 인스턴스 생성에서는 클래스를 호출할 때마다 새 인스턴스가 생성되고 \zsh{__init__()}이 실행된다.

따라서 \zsh{p1}과 \zsh{p2}는 각각 별도의 속성을 갖게 된다.

\begin{dhkcode}
print(p1.__dict__)
print(p2.__dict__)
\end{dhkcode}

실행 결과:

\begin{dhkzsh}
{'name': 'Kim'}
{'name': 'Lee'}
\end{dhkzsh}

여기서 2.3과 2.4에서 학습한 개념들이 연결된다.

\begin{dhksmry}
    - 클래스 객체는 하나다.

    - 인스턴스 객체는 두 개다.

    - 각 인스턴스는 독립적인 속성 저장 공간을 가질 수 있다.

    - __init__()은 각 인스턴스를 초기화하는 과정에서 호출된다.
\end{dhksmry}

다만 앞에서 본 것처럼 \zsh{__new__()}가 다른 객체를 반환하거나 메타클래스의 호출 동작을 변경하면 이러한 일반적인 흐름이 달라질 수 있다.

\subsubsection{PyTorch와의 연결 - \zsh{super().__init__()}은 왜 필요한가?}

이제 PyTorch 모델 코드로 돌아가 보자.

\begin{dhkcode}
import torch.nn as nn

class MyModel(nn.Module):
    def __init__(self):
        super().__init__()

        self.linear = nn.Linear(10,5)

    def forward(self, x):
        return self.linear(x)
\end{dhkcode}

여기서 가장 중요한 부분은 다음 코드다.

\begin{dhkcode}
super().__init__()
\end{dhkcode}

왜 이 코드를 작성할까?

\zsh{MyModel}은 \zsh{nn.Module}을 상속받는 클래스다.

\zsh{nn.Module}은 하위 모듈과 파라미터를 등록하고 관리하기 위한 내부 상태를 가지고 있다.

따라서 \zsh{MyModel}의 초기화 과정에서는 먼저 \zsh{nn.Module}이 필요로 하는 내부 상태를 초기화해야 한다.

그 역할을 수행하는 것이 다음 코드다.

\begin{dhkcode}
super().__init__()
\end{dhkcode}

이 예제에서는 \zsh{super()}가 MRO에 따라 \zsh{nn.Module}의 \zsh{__init__()}을 찾아 호출하게 된다.

이후 다음 코드가 실행된다.

\begin{dhkcode}
self.linear = nn.Linear(10,5)
\end{dhkcode}

이때 \zsh{nn.Module}의 내부 상태가 초기화되어 있으므로, \zsh{nn.Module} 인스턴스를 하위 모듈로 등록할 수 있다.

\begin{dhknote}
\begin{dhksmry}[centered]
초기화 순서가 바뀌면 어떻게 될까?
\end{dhksmry}
\end{dhknote}

\begin{dhkcode}
class WrongModel(nn.Module):
    def __init__(self):
        self.linear = nn.Linear(10,5)

        super().__init__()
\end{dhkcode}

이 코드는 일반적인 PyTorch \zsh{nn.Module}에서 오류를 발생시킨다.

대표적인 오류 메시지는 다음과 같다.

\begin{dhkzsh}
AttributeError: cannot assign module before Module.__init__() call
\end{dhkzsh}

왜 이런 문제가 발생할까?

\zsh{nn.Module}이 하위 모듈을 등록하기 위한 내부 관리 구조를 아직 초기화하지 않았기 때문이다.

따라서 일반적인 PyTorch 모델에서는 다음 순서를 사용해야 한다.

\begin{dhkcode}
class MyModel(nn.Module):
    def __init__(self):
        super().__init__()

        self.linear = nn.Linear(10,5)
\end{dhkcode}

이 과정을 그림으로 나타내면 다음과 같다.

\begin{dhksmry}[centered]
MyModel()
\da
인스턴스 생성
\da
MyModel.__init__()
\da
super().__init__()
\da
nn.Module 내부 상태 초기화
\da
self.linear = nn.Linear(10,5)
\da
하위 모듈 등록
\da
모델 초기화 완료
\end{dhksmry}

여기서 주의할 점은 \zsh{super()}가 무조건 직접적인 부모 클래스만 호출하는 문법은 아니라는 사실이다.

정확히는 현재 클래스와 MRO를 기준으로 다음 메서드를 탐색하는 구조다.

이 부분은 2.9에서 자세히 다룰 것이다.

\subsubsection{핵심 정리}

이번 학습에서는 클래스 객체 호출과 인스턴스 생성 및 초기화의 관계를 살펴보았다.

\begin{dhknote}
\begin{dhksmry}[centered]
__new__()와 __init__()은 서로 다른 역할을 한다.
\end{dhksmry}
\end{dhknote}

\begin{dhkcode}
def __new__(cls, ...):
    ...
\end{dhkcode}

인스턴스 생성에 관여한다.

\begin{dhkcode}
def __init__(self, ...):
    ...
\end{dhkcode}

이미 생성된 인스턴스를 초기화한다.

\begin{dhknote}
\begin{dhksmry}[centered]
일반적인 클래스 호출은 type.__call__()을 통해 처리된다.
\end{dhksmry}
\end{dhknote}

\begin{dhkcode}
p = Person("Kim")
\end{dhkcode}

기본적인 생성 과정에서는 \zsh{__new__()}와 \zsh{__init__()}이 순서대로 호출된다.

\begin{dhknote}
\begin{dhksmry}[centered]
cls와 self는 서로 다른 객체를 참조한다.
\end{dhksmry}
\end{dhknote}

\begin{dhksmry}
    - cls는 클래스 객체를 참조한다.

    - self는 인스턴스 객체를 참조한다.
\end{dhksmry}

\begin{dhknote}
\begin{dhksmry}[centered]
__init__()은 반드시 None을 반환해야 한다.
\end{dhksmry}
\end{dhknote}

\begin{dhkcode}
def __init__(self):
    self.name = "Kim"
\end{dhkcode}

반환값을 명시하지 않는 것이 일반적이다.

\begin{dhknote}
\begin{dhksmry}[centered]
__new__()가 다른 타입의 객체를 반환하면
__init__()이 호출되지 않을 수 있다.
\end{dhksmry}
\end{dhknote}

\begin{dhkcode}
def __new__(cls):
    return "Hello"
\end{dhkcode}

이 경우 일반적인 \zsh{Person} 인스턴스 초기화가 이루어지지 않는다.

\begin{dhknote}
\begin{dhksmry}[centered]
인스턴스 생성과 초기화는 별도의 과정이다.
\end{dhksmry}
\end{dhknote}

\begin{dhkcode}
p = Person.__new__(Person)

Person.__init__(p, "Kim")
\end{dhkcode}

이 코드는 두 과정의 차이를 이해하기 위한 예제이다.

\begin{dhknote}
\begin{dhksmry}[centered]
__new__()는 생성 과정 자체를 제어할 때 사용한다.
\end{dhksmry}
\end{dhknote}

특히 불변 객체의 서브클래싱이나 특수한 객체 생성 전략에서 활용할 수 있다.

\begin{dhknote}
\begin{dhksmry}[centered]
PyTorch에서는 super().__init__()이 중요하다.
\end{dhksmry}
\end{dhknote}

\begin{dhkcode}
class MyModel(nn.Module):
    def __init__(self):
        super().__init__()
\end{dhkcode}

\zsh{nn.Module}의 내부 상태를 초기화해야 하위 모듈과 파라미터를 정상적으로 등록하고 관리할 수 있다.

\subsubsection{2.6에서 반드시 이해해야 할 결론}

이번 학습의 핵심은 다음 구조다.

\begin{dhksmry}[centered]
Class Object Call
\da
__new__
\da
Instance Creation
\da
__init__
\da
Instance Initialization
\da
Return and Name Binding
\end{dhksmry}

단, 이것은 \zsh{__new__()}가 적절한 인스턴스를 반환하고 기본적인 클래스 호출 규칙을 따르는 일반적인 경우의 흐름이다.

이제 다음 코드가 단순히 \zsh{Person}이라는 클래스로 객체를 만든다는 의미를 넘어,

\begin{dhkcode}
p = Person("Kim")
\end{dhkcode}

\begin{dhksmry}
클래스 객체 호출 \ra 인스턴스 생성 \ra 인스턴스 초기화 \ra 반환값 바인딩
\end{dhksmry}

이라는 실행 과정으로 이해되어야 한다.

다음 2.7에서는 지금까지 학습한 인스턴스 네임스페이스와 메서드 바인딩을 기반으로 Python이 같은 이름의 속성을 어떤 순서로 탐색하는지 살펴볼 것이다.

특히 클래스 속성과 인스턴스 속성의 차이, Mutable 클래스 속성의 공유 문제, 속성 탐색 순서, 디스크립터의 우선순위를 연결해서 설명한다.
